% FIGURE NEEDED: spherical-triangle sketch with pole P, points A, B and the
% northernmost point N of the path AB, sides 90-phi_1, 90-phi_2 and angle
% L_2-L_1 at P; 2022 theory solutions, page 6 (problem 7).

In the diagram, $N$ is the northernmost point of the path
$\overset{\frown}{AB}$. Triangle $APB$, triangle $APN$ and triangle $BPN$ are
all spherical triangles. We shall use the convention that North and West are
positive, and South and East are negative. We notice,
\begin{align*}
  \overset{\frown}{PN} &= -\delta_F = 30^\circ 4'\\
  \overset{\frown}{PA} &= 90^\circ - \phi_A = 48^\circ 17'\\
  \overset{\frown}{PB} &= 90^\circ - \phi_B\\
  \sphericalangle PNA &= \sphericalangle PNB = 90^\circ\\
  \sphericalangle BPA &= \lambda_B - \lambda_A
\end{align*}
\hfill(5.0)
\begin{align*}
  x = \overset{\frown}{AB} &= vt = 4375 \times 2\tfrac{17}{60}
    = 9990\ \mathrm{km}\\
  \therefore x &= \frac{9900}{R_\oplus} = 1.5663\ \mathrm{rad}
  % sic: the official solution writes 9900 here; from the previous line it
  % should be 9990 km.
  \\
  &= 89.74^\circ = 89^\circ 44'
\end{align*}
\hfill(2.0)

In $\triangle APN$, using sine rule,
\begin{align*}
  \frac{\sin\sphericalangle A}{\sin\overset{\frown}{PN}}
    &= \frac{\sin\sphericalangle N}{\sin\overset{\frown}{PA}}\\
  \therefore \sin\sphericalangle A
    &= \frac{\sin 90^\circ \sin(30^\circ 4')}{\cos\phi_A}
     = \frac{\sin(30^\circ 4')}{\cos(41^\circ 43')} = 0.67119\\
  \therefore \sphericalangle A
    &= 0.7358\ \mathrm{rad} = 42.16^\circ = 42^\circ 10'
\end{align*}
\hfill(3.0)

(Recognise that to have a northern point on the path, it must be the acute
solution)

In $\triangle PBA$, using cosine rule,
\begin{align*}
  \cos\overset{\frown}{PB}
    &= \cos\overset{\frown}{PA}\cos\overset{\frown}{AB}
     + \sin\overset{\frown}{PA}\sin\overset{\frown}{AB}
       \cos\sphericalangle A\\
  \sin\phi_B &= \sin\phi_A \cos x + \cos\phi_A \sin x \cos A\\
    &= \sin(41^\circ 43')\cos(89^\circ 44')
     + \cos(41^\circ 43')\sin(89^\circ 44')\cos(42^\circ 10')\\
  \sin\phi_B &= 0.5563\\
  \therefore \phi_B &= 0.5900\ \mathrm{rad} = 33.80^\circ = 33^\circ 48'
\end{align*}
\hfill(4.0)

(Only one solution in the valid values of latitude)

Again using cosine rule,
\begin{align*}
  \cos\overset{\frown}{AB}
    &= \cos\overset{\frown}{PA}\cos\overset{\frown}{PB}
     + \sin\overset{\frown}{PA}\sin\overset{\frown}{PB}
       \cos\sphericalangle P\\
  \therefore \cos\sphericalangle P
    &= \frac{\cos x - \sin\phi_A \sin\phi_B}{\cos\phi_A \cos\phi_B}\\
    &= \frac{\cos(89^\circ 44') - \sin(41^\circ 43')\sin(33^\circ 48')}
            {\cos(41^\circ 43')\cos(33^\circ 48')}\\
  \cos\sphericalangle P &= -0.5891\\
  \therefore \sphericalangle P
    &= 2.201\ \mathrm{rad} = 126.13^\circ = 126^\circ 8'\\
  \lambda_B &= \sphericalangle P + \lambda_A = 126^\circ 8' - 41^\circ 48'\\
  \lambda_B &= 1.472\ \mathrm{rad} = 84.33^\circ = 84^\circ 20'
\end{align*}
\hfill(5.0)

(The other solution of $\lambda_B - \lambda_A = -126.13^\circ$ gives
$\lambda_B = -167.93^\circ$ which is clearly not valid)

Therefore, the coordinates of Atlanta are
\[
  (33^\circ 48'\mathrm{N}, 84^\circ 20'\mathrm{W})
\]
\hfill(1.0)

\medskip
\noindent\textbf{Alternative method to find $L_2$}

Using the spherical sine rule in triangle $PBA$,
\begin{align*}
  \frac{\sin A}{\sin(90^\circ - \phi_2)}
    &= \frac{\sin(L_2 - L_1)}{\sin x}\\
  \therefore \sin(L_2 - L_1) &= \frac{\sin x \sin A}{\cos\phi_2}\\
  \therefore \sin(L_2 - L_1)
    &= \frac{\sin(89^\circ 51')\sin(42^\circ 10')}{\cos(33^\circ 43')}
     = 0.80689\\
  \therefore L_2 - L_1
    &= 126.21^\circ \quad (= 126^\circ 12' = 2.20\ \mathrm{rad})\\
  \therefore L_2 &= 126^\circ 12' + L_1
    = 126^\circ 12' + (-41^\circ 48') = 84.41^\circ
    \quad (= 84^\circ 24' = 1.47\ \mathrm{rad})
\end{align*}

(The other solution of $L_2 - L_1 = 53.79^\circ$ gives $L_2 = 11.99^\circ$
which is clearly not valid)

Accept all other valid solutions (such as the four parts rule or using
Napier's rules) that give $\phi_2 = 33.71^\circ$ and $L_2 = 84.41^\circ$.
Students may do both methods for calculating $L_2$ to confirm which solution
is common to both as a quicker check than trying to do the reverse calculation
with a given pair of co-ordinates.
